\chapter{Second-Order Ordinary Differential Equations}

A \textbf{second-order ODE} contains the second derivative $d^2y/dx^2$ as the highest-order term. The general linear second-order ODE is
\[
  a(x)\,\frac{d^2y}{dx^2} + b(x)\,\frac{dy}{dx} + c(x)\,y = f(x).
\]

\begin{example}[Classification]
\[
  \frac{d^2y}{dx^2} + 3\,\frac{dy}{dx} + y = x
\]
\begin{itemize}
  \item \textbf{Order:} second (highest derivative is $d^2y/dx^2$).
  \item \textbf{Linearity:} linear ($y$, $y'$, $y''$ all appear to the first power, no products).
  \item \textbf{Type:} non-homogeneous (right-hand side $f(x) = x \neq 0$).
\end{itemize}
\end{example}

\begin{remark}
The full treatment of second-order ODEs — including the method of undetermined coefficients, variation of parameters, and the Wronskian — is developed in the course on Differential Equations (see the dedicated subject notes). The example above illustrates classification only.
\end{remark}
